% LSPV-5 REPORT::PHCLOSE1::2026-10-02T09:42:00-07:00
\documentclass[11pt]{article}
\usepackage[margin=.88in]{geometry}
\usepackage{amsmath,amssymb}
\begin{document}
\section*{PHCLOSE1: bounded-alternation provenance closure under the counterfactual P=NP}

\textbf{STATUS: PROVED/EXACT closure theorem + COMPUTATIONAL sanity test. P versus NP remains OPEN.}

\section*{CORPUS REVIEWED}
The current GitHub repository tree was enumerated at 275 paths. This run reviewed root
\texttt{main.tex}; \texttt{AGENTS.md}; \texttt{CODEX.md}; the pnp-research skill;
TASK-001 through TASK-007; verified MAIN snapshots through 0858; all materially relevant
post-0858 specialist changes, including LSPV-2 TURANEX1; predecessor LSPV-4 TRACECLOSE1;
the 2026-10-02 independent review; PRIME-LQP-1 and PRIMECRT1; BITRANK1; the
HEC/VSC-HEC/G-PPC-HEC/CRS-PPC/CCSPC/RCCN/Search--G branch; current product-span and
prime-LQP experiments/results; governance notes; and current GROK/CLAUDE material.
No newer LSPV-4 report exists after TRACECLOSE1, so predecessor preflight is nonblocking.

\section*{DUPLICATION / CONTRADICTION AUDIT}
PRIME-FNP-CLOSE1 covered polynomially balanced NP-verifiable replacement realizations.
TRACECLOSE1 covered polynomial-time functionals of a deterministic polynomial execution
transcript. PHCLOSE1 strictly generalizes the former from NP-verifiable relations to every
polynomially balanced relation whose graph lies in PH. It also covers uniqueness/canonicality and
fixed-bounded-alternation audit predicates. It does not contradict the TFPH literature: that
literature studies search at higher PH levels without assuming the global collapse P=NP.

TURANEX1 is independent of this closure. Its generalized-weight/support-graph theorem is adopted
within the structural branch and does not alter the provenance proof.

\section*{Primary-literature audit}
Krentel's OptP work identifies lexicographically extremal satisfying assignments with
function/oracle complexity near FP with an NP oracle. Kleinberg--Korten--Mitropolsky introduced
total search problems above TFNP in the polynomial hierarchy. Fleming--Imrek--Marciot develop the
total function polynomial hierarchy and explicitly emphasize higher-level and unique-solution
search problems. Chandra--Kozen--Stockmeyer prove the decisive alternation boundary:
bounded alternation yields the polynomial hierarchy, whereas alternating polynomial time equals
PSPACE. These results support the exact boundary proved below but do not themselves provide a SAT
lower bound.

\section*{Theorem PH-PROV-CLOSE1}
Assume
\[
P=NP.
\]
Let \(R(x,w)\) be a polynomially balanced relation: for some polynomial \(p\),
\[
R(x,w)\Longrightarrow |w|\le p(|x|).
\]
Assume the graph language
\[
G_R=\{\langle x,w\rangle:R(x,w)\}
\]
belongs to PH. Then there is a deterministic polynomial-time selector \(F\) such that for every
\[
x\in D_R:=\{x:\exists w\,R(x,w)\},
\]
\[
\boxed{R(x,F(x)).}
\]

\paragraph{Proof.}
Use a fixed-length padded witness encoding of length \(p(|x|)\). For a prefix \(u\), define
\[
E_R(x,u)
\iff
\exists v\in\{0,1\}^{p(|x|)-|u|}:R(x,uv).
\]
Because \(G_R\in PH\), adding one polynomially bounded existential block keeps
\(E_R\in PH\). Under \(P=NP\), the polynomial hierarchy collapses to P, so
\[
E_R\in P.
\]
Starting from the empty prefix, test \(E_R(x,u0)\). Append 0 if true and 1 otherwise.
After \(p(|x|)\) polynomial-time prefix tests, the resulting string \(w\) satisfies
\(R(x,w)\). Therefore \(F\in FP\). \(\square\)

\section*{Corollary 1: uniqueness and canonical binding do not escape}
If \(R(x,\cdot)\) has a unique valid witness, PH-PROV-CLOSE1 reconstructs that unique witness.
More generally, the prefix procedure returns the lexicographically least valid witness. Hence a
canonicality requirement of the form
\[
w=\min_{\rm lex}\{z:R(x,z)\}
\]
does not by itself evade polynomial reacquisition under \(P=NP\), provided the underlying
validity relation remains polynomially balanced and in PH.

This sharpens ISOCAN1/PRIME-FNP-CLOSE1: uniqueness is not a surviving prime-provenance hinge merely
because it prevents substitution by an arbitrary witness.

\section*{Corollary 2: fixed-bounded-alternation audit closure}
Suppose a polynomial-length SAME\_REALIZATION transcript \(w\) is declared valid by a relation
with any fixed constant number of polynomially bounded quantifier alternations, for example
\[
R(x,w)
\iff
\forall y_1\exists y_2\forall y_3\;V(x,w,y_1,y_2,y_3),
\]
with polynomial-time \(V\). Then the graph of \(R\) is in PH. Under \(P=NP\), a valid
transcript is therefore selectable in FP.

Consequently, adding a fixed number of canonical challenge/audit layers to a polynomial transcript
does not escape the counterfactual closure. This applies to any CRS-PPC/CCSPC/PRIME-LQP variant
whose complete legal realization can be expressed as one polynomial-length witness with
fixed-bounded-alternation legality.

\section*{Corollary 3: prime annotation closure}
Let \(g(x,w)\) be any polynomial-time computable positional-prime, CRT, residue, divisibility, or
other polynomial-output annotation. Under the hypotheses of PH-PROV-CLOSE1,
\[
x\mapsto g(x,F(x))
\]
is in FP. Therefore replacing an NP-verifiable prime history by a fixed-PH-verifiable prime history
does not produce a non-P-closed resource under the assumed \(P=NP\).

PRIME-FNP-CLOSE1 is the first-level special case.

\section*{The genuine alternation boundary}
The closure above depends critically on membership in PH, meaning a fixed bound on alternations.
Chandra--Kozen--Stockmeyer show
\[
\boxed{AP=PSPACE.}
\]
Thus a requirement with polynomially many genuine existential/universal alternations can leave PH
and enter PSPACE territory. The assumption \(P=NP\) is not known to imply \(P=PSPACE\).

This is the first syntactic zone not eliminated by PH-PROV-CLOSE1, but it is not automatically
useful for P versus NP. A polynomial-depth challenge/response condition that quantifies over all
future challenges is a strategy object, not merely one realized polynomial transcript. A full
strategy may have exponential explicit size. A succinct strategy representation reintroduces a
verification/representation theorem. Most importantly, the project would still need a
model-preserving theorem forcing every hypothetical polynomial SAT decider to realize that
PSPACE-like strategy condition.

Therefore
\[
\boxed{\text{outside PH}\neq\text{forced SAT resource}.}
\]

\section*{COMPUTATIONAL sanity test}
The committed experiment
\[
\texttt{experiments/prime\_lqp/ph\_provenance\_closure\_test.py}
\]
with raw output
\[
\texttt{results/raw/prime\_lqp/2026-10-02\_LSPV5\_PHCLOSE1.json}
\]
uses seed 20261002 and five witness bits. Fifty toy relations were generated in each of three
classes: an existential graph, a \(\Pi_2\)-style graph, and a \(\Pi_3\)-style graph. An exhaustive
finite prefix-extension oracle reconstructed the lexicographically least valid witness in exactly
five prefix decisions in all 150 cases. The induced canonical witness was unique in every case.

This is only a finite verification of the prefix-search mechanism. It does not simulate \(P=NP\)
and provides no asymptotic lower-bound evidence.

\section*{Implications for the causal/provenance branch}
A provenance proposal now faces the following closure ladder.

\begin{enumerate}
\item If it factors through ordinary semantics, provenance is irrelevant.
\item If it is a polynomial-time functional of a deterministic polynomial trace, TRACECLOSE1 applies.
\item If it accepts a polynomially balanced NP-verifiable replacement realization, PRIME-FNP-CLOSE1 applies.
\item If it accepts any polynomially balanced fixed-PH-verifiable realization, PH-PROV-CLOSE1 applies.
\item To escape this ladder syntactically, the requirement must leave fixed PH, for example through
polynomially many alternations, an unbounded predicate, or an exponentially large strategy object.
Such an escape is useful only if an independent solver-independent compiler/necessity theorem forces it.
\end{enumerate}

This substantially narrows CCSPC-A/CRS-B style canonical-transition ideas. Fixed bounded
alternation does not suffice. A true online strategy condition can avoid PH-PROV-CLOSE1, but then
TASK-004/TASK-005 must prove staticization and arbitrary-algorithm transfer without silently turning
ordinary SAT into a PSPACE problem.

\section*{Relativization, algebrization, and hostile controls}
PH-PROV-CLOSE1 relativizes: if \(P^O=NP^O\), then \(PH^O=P^O\), and the same prefix proof works
relative to \(O\). It is therefore a barrier/closure theorem, not the nonrelativizing hinge.
Prime/CRT arithmetic remains explicitly algebraic.

DEAD/COPY/DELAY can alter provenance while preserving semantics. INTERP/RECOMP regenerate
ordinary deterministic histories. Padding and auxiliary variables can alter explicit transcript
length without changing the semantic problem. Horn, dual-Horn, 2-SAT, affine/Tseitin,
bounded-treewidth, renaming, and independent-product controls all remain compatible with the closure.

\section*{CORPUS RECONCILIATION}
\textbf{New verified result:} PH-PROV-CLOSE1 and its uniqueness/fixed-alternation corollaries.

\textbf{Conflict resolved:} canonical provenance is not a generic escape from reacquisition if
canonical validity lies in PH. Unique TFPH search remains a meaningful literature topic without the
assumption \(P=NP\); under that counterfactual, the direct prefix-collapse argument applies.

\textbf{Stale claim removed/downgraded:} any suggestion that a constant number of extra
universal/existential audit layers alone defeats PRIME-FNP-CLOSE1 is retired.

\textbf{Falsifications retained:} raw prime magnitude/factorization, deterministic trace
provenance, generic canonical binding, causal depth alone, and fixed-bounded-alternation transcript
provenance are not unrestricted SAT lower-bound resources.

\textbf{Adopted handoffs:} TURANEX1 within the structural diagnostic branch; MAIN 0858 BITRANK1
characteristic-two target; TRACECLOSE1; PRIMECRT1; current independent review.

\section*{Single next target}
The overall team target remains structural:
\[
\boxed{\operatorname{rank}_{\mathbb F_2}\Phi=\omega(m)}
\]
or a classification of the \(\Theta(m^2)\)-kernel exceptions on one exact source-typed SU/vMS
family, using characteristic-two certificates with MATRIXCERT1, EXCESS/TURANEX1, and COVERBREAK1
controls.

For provenance specifically, perform only one further falsification test: formalize the smallest
source-legal polynomial-alternation SAME\_REALIZATION strategy relation that does not reduce to one
PH-verifiable realized transcript. If no model-preserving compiler forces that strategy on an
arbitrary static SAT solver, downgrade PRIME-LQP/causal provenance to instrumentation and restricted
model work rather than a proof-spine candidate.

\begin{thebibliography}{9}
\bibitem{CKS81}
A. K. Chandra, D. C. Kozen, L. J. Stockmeyer,
"Alternation," Journal of the ACM 28(1), 1981, 114--133.

\bibitem{Krentel88}
M. W. Krentel,
"The Complexity of Optimization Problems,"
Journal of Computer and System Sciences 36(3), 1988, 490--509.

\bibitem{KKMP21}
R. Kleinberg, O. Korten, D. Mitropolsky, C. Papadimitriou,
"Total Functions in the Polynomial Hierarchy,"
ITCS 2021, LIPIcs 185:44.

\bibitem{FIM25}
N. Fleming, D. Imrek, C. Marciot,
"Provably Total Functions in the Polynomial Hierarchy,"
CCC 2025, LIPIcs 339:28.
\end{thebibliography}

\end{document}
% END_OF_VERIFIED_REPORT::LSPV-5::PHCLOSE1
